% =============================================================================
%  CodeTrackr - Resume Project Entry  (ATS-optimised, Overleaf-ready)
% =============================================================================
%
%  WHAT THIS FILE IS
%  ----------------
%  A ready-to-paste "Projects" entry for CodeTrackr, written to (a) parse
%  cleanly through Applicant Tracking Systems and (b) stay 100% truthful to
%  what the code actually does, so nothing here can be picked apart in an
%  interview.
%
%  HOW TO USE
%  ----------------
%  * Preview it:  this file compiles on its own with pdfLaTeX.
%                 Overleaf -> New Project -> Blank Project -> paste -> Recompile.
%  * Drop it in:  copy the block between  "% >>> COPY FROM HERE"  and
%                 "% <<< COPY TO HERE"  into your existing resume's Projects
%                 section (works with any template - it uses only itemize).
%  * Jake's / sb2nov template users: a macro-format version is at the bottom
%                 of this file, after \end{document}.
%
%  REPLACE THE <ANGLE-BRACKET> PLACEHOLDERS:
%     <DATES>            e.g. "Jan 2024 -- Aug 2026"  (changelog spans 2024-2026)
%     <GITHUB URL>       https://github.com/Soham-Official/CodeTrackr
%     <LIVE URL>         https://code-trackr-frontend.vercel.app   (confirm it's up)
%
%  ATS CHECKLIST (why this entry scores well)
%  ----------------
%  [x] Single column, no tables / text boxes / graphics in the entry
%  [x] \usepackage[T1]{fontenc}  -> text stays selectable & extractable (no "fi" ligature loss)
%  [x] Standard section name: "Projects"
%  [x] Real bullet list (itemize), strong action verbs, one idea per bullet
%  [x] Acronyms expanded on first use:  OAuth 2.0, JSON Web Token (JWT), IDOR
%  [x] Tech keywords appear BOTH in the heading and in a Technical Skills block
%  [x] Plain-text dates; nothing important in headers/footers
% =============================================================================

\documentclass[11pt,a4paper]{article}

\usepackage[T1]{fontenc}      % selectable/extractable text for ATS
\usepackage[utf8]{inputenc}
\usepackage{lmodern}          % Latin Modern: correct ToUnicode -> clean ATS text extraction
% Force a ToUnicode map so ATS parsers read "file"/"efficient", not "le"/"ecient".
% (Guarded so the file still compiles under XeLaTeX / LuaLaTeX; use pdfLaTeX for best results.)
\ifdefined\pdfgentounicode
  \input{glyphtounicode}
  \pdfgentounicode=1
\fi
\usepackage[margin=1.6cm]{geometry}
\usepackage{enumitem}
\usepackage[hidelinks]{hyperref}
\usepackage{titlesec}

\titleformat{\section}{\large\bfseries}{}{0pt}{}[\vspace{2pt}\hrule\vspace{4pt}]
\setlist[itemize]{leftmargin=1.4em,itemsep=2pt,topsep=3pt,parsep=0pt}
\pagestyle{empty}
\setlength{\parindent}{0pt}

\newcommand{\project}[4]{%
  % #1 name  #2 tech stack  #3 dates  #4 itemize block
  \textbf{#1} \hfill \textbf{#3}\\[1pt]
  \textit{#2}\\[-6pt]
  #4\vspace{4pt}
}

\begin{document}

% -----------------------------------------------------------------------------
\section*{Projects}
% >>> COPY FROM HERE >>>------------------------------------------------------

\project
  {CodeTrackr: Developer Productivity Analytics Platform}
  {TypeScript, JavaScript, React, Node.js, Express.js, MongoDB, Mongoose,
   REST APIs, OAuth 2.0, JSON Web Token (JWT), VS Code Extension API, Vite, Chart.js, Git}
  {<DATES>}
  {%
  \begin{itemize}
    \item Built a full-stack developer-productivity analytics platform with three
          independent clients served by one REST API: a \textbf{TypeScript}
          \textbf{VS Code extension}, a \textbf{React} / \textbf{Vite} dashboard, and
          a \textbf{Node.js} / \textbf{Express.js} backend on \textbf{MongoDB}.
    \item Engineered the extension's event-driven tracking engine (editor, terminal,
          git, focus, and debug trackers) that accumulates coding metrics in memory
          and flushes normalised activity documents to the API on a 30-second
          interval, transmitting only file basenames and sanitised command metadata
          to keep source code private.
    \item Implemented authentication and authorisation end to end: \textbf{Google
          OAuth 2.0} with \textbf{JWT} session cookies for the web client, per-user
          API keys for the extension, scrypt-hashed passwords for private groups,
          and route-level ownership checks that eliminated an insecure direct object
          reference (IDOR) vulnerability class.
    \item Modelled time-series activity data in \textbf{MongoDB} and authored
          \textbf{aggregation pipelines} plus a unit-tested metrics service that
          derives productivity insights (deep-work ratio, day-to-day consistency via
          coefficient of variation, peak-productivity hour, estimate-vs-actual
          accuracy) from raw events.
    \item Ran a self-directed security and correctness audit that catalogued 33
          findings by severity; remediated every high-severity issue (broken access
          control, plaintext credential storage, an authentication-bypass flag) and
          wrote 8 automated test suites, including a static scanner that fails when a
          protected route loses its authorisation middleware.
    \item Deployed the dashboard on \textbf{Vercel}, the API on \textbf{Render}, and
          the database on \textbf{MongoDB Atlas}; packaged and published the
          extension to the \textbf{VS Code Marketplace}.
    % OPTIONAL extra line: add only if you have room / want the links inline.
    % \item Code: \href{<GITHUB URL>}{github.com/Soham-Official/CodeTrackr}
    %       \quad Live: \href{<LIVE URL>}{code-trackr-frontend.vercel.app}
  \end{itemize}}

% <<< COPY TO HERE <<<-------------------------------------------------------

% -----------------------------------------------------------------------------
\section*{Technical Skills}
% Repeat the CodeTrackr keywords here too -- ATS weight terms that appear in a
% dedicated skills section. Merge these into your existing lines.

\textbf{Languages:} TypeScript, JavaScript, Python, SQL, HTML/CSS\\
\textbf{Frameworks \& Libraries:} React, Node.js, Express.js, Mongoose, Vite, Chart.js, Tailwind CSS\\
\textbf{Databases:} MongoDB, MongoDB Atlas, PostgreSQL\\
\textbf{Concepts:} REST API design, Authentication \& Authorisation, OAuth 2.0, JWT,
  Data Aggregation Pipelines, Time-Series Data, Unit Testing, Git, Cloud Deployment\\
\textbf{Tools \& Platforms:} Git, GitHub, Vercel, Render, VS Code Extension API, Postman

\end{document}

% #############################################################################
% #  ALTERNATE FORMATS + TAILORING NOTES  (LaTeX ignores everything below      #
% #  \end{document}, so this is just reference text -- copy what you need.)    #
% #############################################################################


% =============================================================================
%  OPTION B - "Jake Gutierrez" resume template  (\resumeProjectHeading macros)
%  Also works with sb2nov. Paste inside \resumeSubHeadingListStart ... End.
% =============================================================================
%
% \resumeProjectHeading
%   {\textbf{CodeTrackr} $|$ \emph{TypeScript, React, Node.js, Express.js, MongoDB, OAuth 2.0, JWT, VS Code Extension API}}
%   {<DATES>}
%   \resumeItemListStart
%     \resumeItem{Built a full-stack developer-productivity analytics platform with three
%       clients on one REST API: a TypeScript VS Code extension, a React/Vite dashboard,
%       and a Node.js/Express.js backend on MongoDB.}
%     \resumeItem{Engineered the extension's event-driven tracking engine (editor, terminal,
%       git, focus, debug trackers) that flushes normalised activity documents to the API
%       every 30 seconds, sending only file basenames and sanitised command metadata.}
%     \resumeItem{Implemented Google OAuth 2.0 with JWT session cookies, per-user API keys
%       for the extension, scrypt-hashed group passwords, and route-level ownership checks
%       that removed an IDOR vulnerability class.}
%     \resumeItem{Wrote MongoDB aggregation pipelines and a unit-tested metrics service
%       deriving productivity insights (deep-work ratio, consistency index, peak hour)
%       from time-series activity data.}
%     \resumeItem{Ran a 33-finding security/correctness audit; fixed all high-severity
%       issues (broken access control, plaintext credentials, auth-bypass flag) and added
%       8 automated test suites incl. a static route-guard scanner.}
%   \resumeItemListEnd


% =============================================================================
%  COMPACT VERSION - 2 lines, for a dense one-page resume
% =============================================================================
%
% \textbf{CodeTrackr} $|$ React, TypeScript, Node.js, Express.js, MongoDB, OAuth 2.0, JWT, VS Code Extension API \hfill <DATES>
% \\ Full-stack coding-analytics platform (VS Code extension + React dashboard + Express/MongoDB REST API): an
% event-driven extension streams activity to the API, which surfaces productivity metrics via MongoDB aggregation
% pipelines; ran a 33-item security audit and fixed all high-severity access-control and credential-storage issues,
% backed by an 8-suite automated test set.


% =============================================================================
%  KEYWORD TAILORING - edit per job description before you apply
% =============================================================================
%  * Mirror the JD's exact words:
%       "REST API"        -> "RESTful API"        if the JD says RESTful
%       "authorisation"   -> "authorization"      match the JD's spelling (US vs UK)
%       "unit testing"    -> "automated testing"  / "test-driven development" if claimed
%  * Safe additions if the JD lists them AND you can defend them:
%       "Agile", "code review", "responsive design", "SPA", "CI/CD" (only after you
%       add a GitHub Actions workflow), "npm", "ESLint", "TypeScript strict mode".
%  * DO NOT add: "microservices" (it is a monolith), "Docker/Kubernetes" (none),
%       "GraphQL" (none), "Redis/Kafka" (none), "machine learning" (see note below),
%       "10k+ users" / any usage or performance number (never load-tested).


% =============================================================================
%  INTERVIEW-SAFETY NOTES - so these bullets never backfire
% =============================================================================
%  * "Metrics service / productivity insights": say STATISTICAL ANALYSIS, not
%    machine learning. It is coefficient of variation, weighted scoring and
%    medians - no model, no training. This is a strength (every number is
%    explainable); own it as a deliberate choice.
%  * The security-audit bullet is your strongest differentiator. Be ready to:
%      - name the fixes: analytics endpoints had no ownership check (IDOR);
%        private-group passwords were stored in plain text; an AUTH_BYPASS env
%        flag disabled all auth; two legacy endpoints accepted unauthenticated writes.
%      - define IDOR in one line: "any logged-in user could read another user's
%        data by putting their id in the URL - I added a check that the record
%        belongs to the caller."
%  * "Flushes every 30 seconds": the timer ticks every 30s; a network send
%    happens once enough *active* coding time has accrued (default ~30s). Know
%    the nuance.
%  * "Published to the VS Code Marketplace": earlier versions were published
%    under publisher "CodeTrackr-ext"; confirm the current build is live before
%    you say "currently published".
%  * If asked how it scales: it is a portfolio project, not load-tested. Talk
%    about the plan - a per-user UserStats rollup so the leaderboard is an
%    indexed 50-row read instead of a full-collection scan, and Redis for rank
%    lookups. That answer shows more than a fake metric would.
%  * Only claim first-person work you actually did (repo has 3 contributors).
%    If the audit + tests were yours, keep that bullet; if not, swap it for the
%    deployment bullet.
% =============================================================================
